\section{楚：吴起改革与南方大国的内部阻力}

吴起先后活动于魏、楚，正好让我们看见人才和制度经验怎样跨国流动。传统材料将其在楚的改革与削弱贵族封君、整顿吏治联系在一起，也记载其死后改革受反弹。

楚的问题不是“南方所以落后”，而是疆域辽阔、旧贵族力量强、中央改革成本高。吴起的遭遇说明：一项改革能否留下来，不只取决于条文是否有效，还取决于改革者是否能重组支持者、继承安排与利益补偿。

\begin{sourcenote}[吴起故事怎样使用]
\sourcebook{史记·孙子吴起列传}是吴起最系统的传记，写在其身后数百年；《战国策》的说辞和《吴子》的篇章也都有后人整理、文学化或复合成书的问题。吴起在魏训练武卒、在楚削弱世卿并在楚悼王死后被贵族杀害，是可以谨慎把握的长线；具体选兵标准和改革条文则不宜写得像现代档案。
\end{sourcenote}
